%% kapa: exam=1 round=35 year=2024 period=1 form=A subject=civil_law subject_ko=민법 seq=35 number=35
%% source: Q-Net 감정평가사 기출문제 (cst00309) artlSeq=5213229 file=raw/1008/5213229/1교시 문제지 원본.pdf
%% conversion: kordoc 4.24.51 -> tools/gen1.py (auto); review=proofread-sonnet; flags=none
\begin{question}{35}
유치권에 관한 설명으로 옳은 것은? (다툼이 있으면 판례에 따름)
\begin{choices}
  \item 피담보채권이 존재한다면 타인의 물건에 대한 점유가 불법행위로 인한 것인 때에도 유치권이 성립한다.
  \item 유치권자가 유치물 소유자의 승낙 없이 유치물을 임대한 경우, 특별한 사정이 없는 한 유치물의 소유자는 유치권의 소멸을 청구할 수 없다.
  \item 목적물에 대한 점유를 상실한 경우, 유치권자가 점유회수의 소를 제기하여 점유를 회복할 수 있다는 것만으로는 유치권이 인정되지 않는다.
  \item 채무자를 직접점유자로 하여 채권자가 간접점유를 하였더라도 채권자는 유효하게 유치권을 취득할 수 있다.
  \item 저당물의 제3취득자가 저당물의 개량을 위하여 유익비를 지출한 때에는 민법 제367조에 의한 비용상환청구권을 피담보채권으로 삼아 유치권을 행사할 수 있다.
\end{choices}
\end{question}
\answer{3}
